\documentclass[border=10pt,crop=true]{standalone}
\usepackage{circuitikz}
\usetikzlibrary{positioning}

\begin{document}
\begin{circuitikz}[american, scale=1.2, transform shape, font=\small]
  \coordinate (P) at (0, 0);
  \coordinate (N) at (0, -2.6);

  % terminals
  \draw (P) node[left] {p} to[short, o-] ++(0.4, 0) coordinate (Ptop);
  \draw (N) node[left] {n} to[short, o-] ++(0.4, 0) coordinate (Nbot);

  % diode branch (left)
  \draw (Ptop) to[D, l=$D$, *-] (Nbot);

  % photocurrent source branch (right, parallel)
  \draw (Ptop) -- ++(1.6, 0) coordinate (Pbr)
    to[I, l=$I_{\mathrm{ph}}$, v=$I_{\mathrm{ph}}$] ++(0, -2.6)
    -- (Nbot);

  % terminal voltage v
  \draw[<-, thick] ($(P)+(-0.15,-0.55)$) -- node[left] {$+$} ($(P)+(-0.15,-1.05)$);
  \draw[->, thick] ($(N)+(-0.15,0.55)$) -- node[left] {$-$} ($(N)+(-0.15,1.05)$);
  \node at (-0.15, -0.05) {$v$};

  % total terminal current i (p to n)
  \draw[->, thick] (2.35, -3.15) -- node[above, font=\scriptsize] {$i$} (3.55, -3.15);

\end{circuitikz}
\end{document}
